\begin{textsample}{POS Dim 1 – rr – Score 67.00 – rr/rr\_em\_36.txt}  \label{ex:f1_pos_211}
4.5.
Área Ciências da Natureza O documento apresenta referências para a implantação do currículo da área Ciências da Natureza para o Ensino Médio no âmbito da BNCC ( BRASIL, 2018a ).
Este documento de caráter normativo consolida o compromisso com o \textbf{letramento} científico do estudante, a fim de que seja capaz de compreender e interpretar o mundo ( natural, social e tecnológico ), e consiga ainda desenvolver as habilidades e \textbf{competências} necessárias para transformá -lo.
O Ensino de Ciências, de uma \textbf{maneira} geral, tem reforçado a visão de \textbf{ciência} como algo estático, um conjunto de verdades imutáveis, de estruturas \textbf{conceituais} congeladas no tempo.
Muitas vezes, não tem nenhuma relação com os contextos históricos, sociais e tecnológicos em que a \textbf{ciência} é construída e aplicada.
A ausência de fenômenos não favorece a percepção da natureza das construções \textbf{teóricas} e dos modelos científicos, como construções matemáticas e \textbf{discursivas} para interpretação e \textbf{descrição} de uma realidade muito mais \textbf{complexa} ( \textbf{QUADROS}, MORTIMER, 2016 ).
Conforme Quadros e Mortimer ( 2016 ), a ausência de diálogo entre a realidade criada pela \textbf{ciência} e a realidade da vida cotidiana, entre a linguagem científica e a cotidiana, não possibilita ao jovem rever seu conhecimento à luz das \textbf{novidades} que aprende nas \textbf{aulas} das \textbf{disciplinas} que compõem a área de \textbf{Ciências} da Natureza, fazendo com que não ocorra um diálogo entre as \textbf{teorias} científicas e os fenômenos, entre os princípios científicos e os contextos sociais e tecnológicos em que eles se \textbf{materializam}.
O ensino de Ciências da Natureza no Ensino Fundamental tem como intencionalidade cooperar para a transformação da sociedade, ao tratar dos saberes que lhes são \textbf{inerentes}, permitindo ao estudante o desenvolvimento de habilidades para a construção, reconstrução ou desconstrução dos conhecimentos, premissa que requer a implementação de um conjunto de \textbf{encaminhamentos} que contribuam para a formação de estudantes questionadores e investigativos.
Essa \textbf{capacidade} investigativa dos estudantes, por sua vez, se dá por meio da observação e da pesquisa com o objetivo da integração do conhecimento científico \textbf{resultante} da investigação da natureza para interpretar racionalmente os fenômenos naturais observados, decorrentes de contextos históricos, sociais e \textbf{econômicos} considerados no tempo, espaço, matéria, movimento, força, campo, energia, vida e evolução.
Sendo assim, o ensino de Ciências no Ensino Fundamental deve promover situações nas quais os \textbf{alunos} possam : observar, analisar, propor, planejar, investigar, relatar, desenvolver e implementar ações de intervenção para melhorar a qualidade de vida individual, coletiva e \textbf{socioambiental} ( BRASIL, 2018a ).
Essas diversas situações corroboram para um conjunto de habilidades que também são propostas pela BNCC ( BRASIL, 2018a ), quando \textbf{sugere} uma mudança de paradigma, na qual estrutura o conjunto de habilidades cuja complexidade cresce progressivamente ao longo dos anos.
Resultando numa distribuição de conteúdos tradicionais do componente curricular mais equilibrada.
Antes, o \textbf{foco} em Biologia era maior, com Física e Química sendo abordados com maior frequência, apenas nos anos \textbf{finais} do Ensino Fundamental.
Agora essas áreas das Ciências estão distribuídas nas unidades temáticas e são trabalhadas em todos os anos da escolaridade.
Ou seja, ocorre um trabalho em espiral, em que os eixos se repetem a cada ano, com a \textbf{indicação} de uma \textbf{progressão} de aprendizagem, com os conceitos sendo construídos gradativamente, com complexidade maior a cada ano, conforme avança o desenvolvimento e a maturidade dos estudantes ( BRASIL, 2018a ).
Assim, o estudante que, antes, tinha contato com os objetos de conhecimento relacionados à unidade temática Matéria e Energia apenas no 5º ano do Ensino Fundamental, para depois voltar a estudá -los somente no 9º ano, agora aprenderá as noções mais básicas da área desde os primeiros anos do Ensino Fundamental.
A expectativa é que, quando as fórmulas e cálculos forem apresentados no Ensino Médio, o estudante esteja familiarizado com o \textbf{fundamento} desses objetos de conhecimento.
No Ensino Médio, proposto pela BNCC ( BRASIL, 2018a ), a área de \textbf{Ciências} da Natureza deve se comprometer com a formação dos jovens para o enfrentamento dos desafios da \textbf{contemporaneidade}, na direção da educação integral e da formação cidadã.
No Ensino Fundamental foi enfatizado o \textbf{letramento} científico como compromisso da área envolvendo a \textbf{capacidade} de compreender e interpretar o mundo, bem como transformá -lo com base em aportes \textbf{teóricos} e processuais das \textbf{ciências}, assegurando aos estudantes dessa etapa o acesso a diversos conhecimentos, os \textbf{aproximando} paulatinamente aos principais processos, práticas e procedimentos da investigação científica ( BRASIL, 2018a ).
É necessário destacar que, em especial nos dois primeiros anos da escolaridade básica, em que se \textbf{investe} prioritariamente no processo de alfabetização das crianças, as habilidades de \textbf{Ciências} da Natureza buscam a ampliação dos contextos de \textbf{letramento}.
Nesse sentido, o \textbf{letramento} científico também é enfatizado no Ensino Médio, propiciando uma \textbf{inter-relação} entre o Ensino Fundamental nas habilidades de \textbf{Ciências} da Natureza ( BRASIL, 2018a ).
Para o Ensino Médio, a área das Ciências da Natureza e suas Tecnologias – CNT deve possibilitar o \textbf{aprofundamento} e ampliação dos conhecimentos explorados na etapa anterior, permitindo que os estudantes analisem fenômenos e processos, utilizando modelos e fazendo previsões.
Dessa forma, propicia a \textbf{compreensão} sobre a vida, o nosso planeta e o universo, promovendo a \textbf{capacidade} de refletir, \textbf{argumentar}, propor soluções e enfrentar desafios pessoais e coletivos, locais e globais ( BRASIL, 2018a ).
Essa \textbf{compreensão} sobre a vida \textbf{mobiliza} os estudantes, com maior vivência e maturidade, no \textbf{aprofundamento} do pensamento \textbf{crítico}, onde ele possa realizar novas leituras do mundo, com base em modelos abstratos, e tomar decisões responsáveis, éticas e consistentes na identificação e solução de situações-problema ( BRASIL, 2018a ).
De acordo com Smole e Diniz ( 2001 ), esta solução de situações-problema é vista como perspectiva \textbf{metodológica} de \textbf{ensino-aprendizagem}.
A BNCC da área de CNT, por meio de um olhar \textbf{articulado}, é composta pelos componentes Biologia, Física e Química, com \textbf{competências} e habilidades que permitem a ampliação e a \textbf{sistematização} das aprendizagens essenciais.
Essas aprendizagens são desenvolvidas desde o Ensino Fundamental no que se refere aos conhecimentos \textbf{conceituais} da área ; à \textbf{contextualização} social, cultural, ambiental e histórica desses conhecimentos ; aos processos e práticas de investigação e às linguagens das Ciências da Natureza ( BRASIL, 2018a ).
Sendo assim, podemos considerar importantes no ensino das CNT a \textbf{contextualização}, \textbf{interdisciplinaridade}, produção compartilhada, diversidade cultural, \textbf{problematização}, práticas experimentais, pesquisa de campo e bibliográfica, produção e utilização de texto.
Para isso, é necessário que as escolas desenvolvam diferentes saberes, manifestações culturais e visões de mundo.
Essa instituição deve se constituir como um espaço de heterogeneidade e \textbf{pluralidade}, que valoriza a diversidade e se pauta em princípios de solidariedade e emancipação ( BRASIL, 2018a ).
O ensino de \textbf{Ciências} da Natureza, nos anos iniciais de escolaridade, contribui com a alfabetização, ao mesmo tempo em que proporciona a elaboração de novos conceitos, como propõe Ausubel ( 1963 ) na Teoria de Aprendizagem Significativa, onde o conhecimento do sujeito permite dar significado a um novo conhecimento, seja de forma mediada ou por inferência.
Por isso, a importância do indivíduo trazer para a escola suas vivências e seus saberes, que devem ser tratados de acordo com o que cabe a essa etapa da aprendizagem.
Nos anos \textbf{finais} do Ensino Fundamental, ampliam -se os interesses pela vida social, há uma maior autonomia intelectual.
Fazendo com que se permita o tratamento de conceitos mais amplos e abstratos, que \textbf{dizem} respeito às relações dos sujeitos com a natureza, com as \textbf{tecnologias} e com o ambiente, no sentido da construção de uma visão própria do mundo.
No Ensino Médio, com os jovens e adultos, mais amadurecidos, os conceitos de cada componente curricular — Biologia, Física e Química — podem ser aprofundados em suas especificidades temáticas e em seus modelos abstratos, ampliando a leitura de mundo físico e social, o enfrentamento de situações relacionadas às Ciências da Natureza, o desenvolvimento do pensamento \textbf{crítico} e tomadas de decisões mais conscientes e consistentes.
Para essa formação ampla, os componentes curriculares da área de \textbf{Ciências} da Natureza devem possibilitar a construção de uma base de conhecimentos contextualizada, envolvendo as discussões de temas como energia, saúde, ambiente, \textbf{tecnologia}, educação para o consumo, sustentabilidade, entre outros.
Para que se cumpra com o objetivo de ensino da área de \textbf{Ciências} da Natureza, faz -se necessário integrar os conhecimentos abordados nos componentes curriculares, superando o tratamento de conteúdos fragmentados e articulando com os saberes de outras áreas de conhecimento.
Por exemplo, ao tratar o tema ENERGIA no Ensino Médio, os estudantes, além de compreenderem sua transformação e conservação, do ponto de \textbf{vista} da Biologia, Física e Química, podem percebê -lo sob a perspectiva da Geografia, sabendo avaliar o peso das diferentes fontes de energia na \textbf{matriz} energética.
Considerando fatores como a produção, os recursos naturais \textbf{mobilizados} na produção de energia, as \textbf{tecnologias} envolvidas e os impactos ambientais causados pela produção de energia.
Ainda podemos abordar a apropriação humana dos ciclos energéticos naturais, como elemento essencial para se compreender as transformações \textbf{econômicas} ao longo da História.
Considerando as diversas dimensões formativas citadas, é proposta uma organização dos conhecimentos das CNT em três eixos que possam possibilitar a articulação entre as áreas de conhecimentos ( BRASIL, 2018a ).
Os eixos foram construídos considerando -se as \textbf{competências} específicas do componente curricular de \textbf{Ciências} da Natureza.
No eixo \textbf{Contextualização} social, histórica e cultural das CNT são tratadas as relações entre conteúdos \textbf{conceituais} das Ciências da Natureza e o desenvolvimento histórico da \textbf{ciência} e da \textbf{tecnologia} ; o papel dos conhecimentos científicos e tecnológicos na organização social e formação cultural dos sujeitos e as relações entre \textbf{ciência}, \textbf{tecnologia} e sociedade.
Dessa forma, o currículo deve apontar para estudos de temas de relevância social, a partir dos quais propor articulações entre diferentes áreas ( BRASIL, 2018a ).
Essa \textbf{contextualização} deve contribuir para a aquisição de conhecimentos científicos que possam ser aplicados na vida individual, nos projetos de vida, no mundo de trabalho, vindo a favorecer o protagonismo do estudante no enfrentamento de situações sobre consumo, energia, segurança, ambiente, saúde, entre outras ( BRASIL, 2018a, p. 549 ).
Para isso, as \textbf{competências} específicas e habilidades propostas para o Ensino Médio exploram situações-problema envolvendo melhoria da qualidade de vida, segurança, sustentabilidade, diversidade étnica e cultural, entre outras ( BRASIL, 2018a ).
Em Processos e práticas de investigação em \textbf{Ciências} da Natureza é enfatizada a dimensão do saber fazer, proporcionando ao estudante uma aproximação com os modos de produção do conhecimento científico.
O saber fazer, compreendido não somente como uma metodologia, busca a apropriação desta como um objeto de estudo.
Nesse sentido, o currículo propõe estudos sobre processos de construção de modelos científicos, práticas de investigação científica ( questões e procedimentos de pesquisa adequados ao contexto escolar ), uso e produção de \textbf{tecnologias}, considerando as especificidades do contexto ( BRASIL, 2018a ).
Em Linguagens das Ciências da Natureza é ressaltada a importância de aprender tais linguagens, promovendo a \textbf{compreensão} e a apropriação desse modo de “ se expressar ” próprias da área de CNT.
Destaca -se que essa aprendizagem deve se dar por meio de códigos, símbolos e nomenclaturas próprias da área.
O estudante também deve ter domínio de outras linguagens na comunicação e divulgação do conhecimento científico, reportando -se ao processo de \textbf{letramento} científico necessário a todo cidadão ( BRASIL, 2018a, p. 551 ).
Ainda em relação às Linguagens Específicas da área das CNT, considera -se \textbf{relevante} ressaltar a necessidade de inclusão dos currículos componentes para além do \textbf{aprofundamento} dessas temáticas, ampliando os conhecimentos explorados na etapa anterior.
A BNCC de CNT propõe também que os estudantes ampliem as habilidades investigativas desenvolvidas no Ensino Fundamental, apoiando -se em \textbf{análises} quantitativas e na avaliação e na comparação de modelos explicativos.
Além disso, \textbf{espera} -se que eles aprendam a estruturar linguagens argumentativas que lhes permitam comunicar, para diversos públicos, em contextos variados e utilizando diferentes mídias e \textbf{tecnologias} \textbf{digitais} de informação e comunicação ( TDIC ), conhecimentos produzidos e propostas de intervenção pautadas em evidências, conhecimentos científicos e princípios éticos e responsáveis ( BRASIL, 2018a ).
Assim, a área das CNT coloca a investigação como uma \textbf{maneira} de envolver os estudantes na aprendizagem de processos, práticas e procedimentos científicos e tecnológicos para promover o domínio de linguagens específicas, permitindo a \textbf{análise} de fenômenos e processos por meio de modelos e fazendo previsões ( BRASIL, 2018a, p. 472 ).
Diante disso, neste documento foram \textbf{sugeridas} orientações \textbf{metodológicas} voltadas para atividades práticas contextualizadas e investigação científica, para serem desenvolvidas em consonância com as situações didáticas planejadas pelo professor, que mediará o processo ensino-aprendizagem considerando as quatro modalidades de ação nas quais, de acordo com a BNCC, devem ser : - Definição de \textbf{problemas} : observar o mundo a sua volta e fazer perguntas ; analisar \textbf{demandas}, delinear \textbf{problemas} e planejar investigações ; propor hipóteses. - Levantamento, \textbf{análise} e representação : planejar e realizar atividades de campo ( experimentos, observações, leituras, visitas, ambientes \textbf{virtuais} etc. ) ; desenvolver e utilizar ferramentas, inclusive \textbf{digitais}, para coleta, \textbf{análise} e representação de dados ( imagens, esquemas, tabelas, gráficos, \textbf{quadros}, diagramas, mapas, modelos, representações de sistemas, fluxogramas, mapas \textbf{conceituais}, \textbf{simulações}, \textbf{aplicativos}, etc. ) ; avaliar informação ( validade, coerência e adequação ao \textbf{problema} formulado ) ; elaborar explicações e/ou modelos ; associar explicações e/ou modelos à evolução histórica dos conhecimentos científicos envolvidos ; selecionar e construir \textbf{argumentos} com base em evidências, modelos e/ou conhecimentos científicos ; \textbf{aprimorar} seus saberes e incorporar, gradualmente, e de modo significativo, o conhecimento científico ; desenvolver soluções para \textbf{problemas} cotidianos usando diferentes ferramentas, inclusive \textbf{digitais}. - Comunicação : organizar e/ou extrapolar conclusões ; relatar informações de forma oral, escrita ou multimodal ; apresentar, de forma sistemática, dados e resultados de investigações ; participar de discussões de caráter científico com colegas, professores, familiares e comunidade em geral ; considerar contra-argumentos para rever processos investigativos e conclusões. - Intervenção : implementar soluções e avaliar sua eficácia para resolver \textbf{problemas} cotidianos ; desenvolver ações de intervenção para melhorar a qualidade de vida individual, coletiva e \textbf{socioambiental} ( BRASIL, 2018a, p. 323 ).
De acordo com Santos ( 2007 ), a base para a área de \textbf{Ciências} da Natureza é o \textbf{letramento} científico, definindo que a \textbf{ciência} deve ser usada como ferramenta de atuação sobre o mundo.
O \textbf{letramento} científico consiste na formação técnica do domínio das linguagens e ferramentas mentais usadas em \textbf{ciência} para o desenvolvimento científico.
Nesse sentido, os estudantes deveriam ter amplo conhecimento das \textbf{teorias} científicas e ser capazes de propor modelos em \textbf{ciência}.
Isso \textbf{exige} não só o domínio vocabular, mas a \textbf{compreensão} de seu significado \textbf{conceitual} e o desenvolvimento de processos cognitivos de alto nível de elaboração mental de modelos explicativos para processos e fenômenos.
O \textbf{letramento} científico significa também a busca por uma educação científica que propicie a educação tecnológica.
A \textbf{tecnologia} existe para proporcionar a melhoria da qualidade de vida, tendo como base o consumo consciente e correto de todos os seus produtos que fazem parte do cotidiano.
É necessário compreender, além dos conhecimentos técnico-científicos neles envolvidos, os aspectos éticos e sociais relacionados com a sua produção, comercialização e utilização ( BRASIL, 2006 ).
Na BNCC do Ensino Médio, as aprendizagens essenciais estão organizadas por áreas de conhecimento e para estas são definidas \textbf{competências} específicas, articuladas às respectivas \textbf{competências} das áreas do Ensino Fundamental, sendo adequadas para atender às especificações de formação dos estudantes do Ensino Médio, como também orientar a proposição e o detalhamento dos itinerários formativos relativos às áreas ( BRASIL, 2018a ).
Nesse sentido, para a etapa do Ensino Médio, a BNCC enfatiza que a área de \textbf{Ciências} da Natureza faça a ampliação e \textbf{aprofundamento} das aprendizagens essenciais desenvolvidas durante o Ensino Fundamental no que se refere aos conhecimentos \textbf{conceituais} da área ; \textbf{contextualização} social, cultural, ambiental e histórica ; procedimentos e práticas de investigação e também à linguagem própria das Ciências da Natureza ( BRASIL, 2018a ).
Assim, na definição das \textbf{competências} específicas e habilidades da área de CNT foram privilegiados conhecimentos \textbf{conceituais} dando continuidade à proposta do Ensino Fundamental, considerando sua relevância no ensino de Física, Química e Biologia e sua adequação ao Ensino Médio.
Dessa forma, a BNCC da área de CNT propõe um \textbf{aprofundamento} em três unidades temáticas : Matéria e Energia, Vida e Evolução e Terra e Universo.
A energia está presente em todos os aspectos da vida e, como apresentam Lisboa et al. ( 2016, p. 13 ), “ O ser humano, no \textbf{decorrer} do desenvolvimento das sociedades, aprendeu a utilizar as transformações da matéria para obter energia em diversas formas, entre elas a térmica, elétrica e a luminosa ”, e dentre todas definições de energia pode -se abstrair que energia é a \textbf{capacidade} de realização de um trabalho, movimento ou ação.
Nesse contexto, matéria e energia estão relacionadas, pois todo objeto contém algum tipo de energia.
Para que ocorra alguma transformação da matéria é necessário que ela ganhe ou \textbf{perca} energia.
Em Matéria e Energia, no Ensino Médio, diversificam -se as situações-problema referidas nas \textbf{competências} específicas e nas habilidades, incluindo -se aquelas que permitem a aplicação de modelos com maior nível de abstração e que buscam explicar, analisar e prever os efeitos das interações e relações entre matéria e energia ( por exemplo, analisar \textbf{matrizes} energéticas ou realizar previsões sobre a condutibilidade elétrica e térmica de materiais, sobre o comportamento dos elétrons frente à absorção de energia luminosa, sobre o comportamento dos gases frente a alterações de pressão ou temperatura, ou ainda sobre as consequências de emissões radioativas no ambiente e na saúde ).
Em Vida, Terra e Cosmos, resultado da articulação das unidades temáticas Vida e Evolução e Terra e Universo desenvolvidas no Ensino Fundamental, propõe -se que os estudantes analisem a complexidade dos processos relativos à origem e evolução da Vida ( em particular dos seres humanos ), do planeta, das estrelas e do Cosmos, bem como a dinâmica das suas interações, e a diversidade dos seres vivos e sua relação com o ambiente ( BRASIL, 2018a ).
No entanto, faz -se necessário que o professor busque a prática pedagógica que possibilite a eficácia do processo ensino-aprendizagem, visando potencializar o fazer pedagógico.
Dessa forma, foram propostas orientações didático-metodológicas no \textbf{organizador} curricular como suporte para o desenvolvimento do trabalho docente.
Várias sugestões de estratégias são disponibilizadas como forma de orientar a prática e não como “ receitas ” prontas, já que, conforme Costa e Santos ( 2013 ), a prática docente deve estar atrelada a uma formação contextualizada de acordo com a realidade dos \textbf{discentes} que são público-alvo da mesma, sendo pautada nas perspectivas dos \textbf{alunos}, na tentativa de considerar as particularidades e as \textbf{pluralidades} de cada um ( COSTA ; SANTOS, 2013 ).

% matched lemmas: aluno, análise, aplicativo, aprimorar, aprofundamento, aproximar, argumentar, argumento, articulado, aula, capacidade, ciência, competência, complexo, compreensão, conceitual, contemporaneidade, contextualização, crítico, decorrer, demanda, descrição, digital, discente, disciplina, discursivo, dizer, econômico, encaminhamento, ensino-aprendizagem, esperar, exigir, final, foco, fundamento, indicação, inerente, inter-relação, interdisciplinaridade, investir, letramento, maneira, materializar, matriz, metodológico, mobilizar, novidade, organizador, perder, pluralidade, problema, problematização, progressão, quadro, relevante, resultante, simulação, sistematização, socioambiental, sugerir, tecnologia, teoria, teórico, virtual, vista
\end{textsample}
